\documentclass{article}
\usepackage[utf8]{inputenc}
\title{Pedestrian Crossing Behavior Review}
\begin{document}
\maketitle
\section{Introduction and Research Context}
This review synthesizes evidence from 3 papers examining Predict probability of pedestrian red-light running behavior, Identify significant factors affecting red-light running, Analyze pedestrian red light running decision-making, with findings organized around 8 discrete claims drawn from a structured comparison matrix. The reviewed evidence employs Binomial logistic regression / Statistical modeling, Random parameter logit model (binary logit regression with normally distributed random parameters), Virtual Reality Experiment / Human-subjects laboratory experiment as primary analytical frameworks. A persistent challenge in this literature is the heterogeneity in how studies operationalize crossing behavior constructs, report metrics, and handle context specificity — factors that collectively limit direct cross-study comparison. This review structures its analysis around methodological families, task-level performance patterns, metric reporting practices, and verified gaps (1 identified), moving beyond siloed paper summaries toward a structured evidence synthesis. The review is organized as follows: Section 2 examines methodological approaches and their representational tradeoffs; Section 3 addresses task-level findings and performance patterns; Section 4 evaluates metric reporting and comparability; Section 5 identifies verified research gaps and implications; Section 6 provides a field-level synthesis and research priorities.

\section{Methodological Approaches in Pedestrian Crossing Research}
The Binomial logistic regression / Statistical modeling family constitutes a primary methodological strand in the reviewed literature. Represents a distinct representational approach with characteristic strengths and limitations. Evidence across reviewed studies suggests that this approach is particularly well-suited to analyzing crossing decision determinants under specific traffic conditions. Limitations of this approach include reduced generalizability to contexts with substantially different traffic signal configurations. The Virtual Reality Experiment / Human-subjects laboratory experiment family represents a primary methodological strand in the reviewed literature. Represents a distinct representational approach with characteristic strengths and limitations. Evidence across reviewed studies suggests that this approach is particularly well-suited to analyzing crossing decision determinants under specific traffic conditions. Limitations of this approach include reduced generalizability to contexts with substantially different traffic signal configurations. The Random parameter logit model (binary logit regression with normally distributed random parameters) family represents a secondary methodological strand in the reviewed literature. Represents a distinct representational approach with characteristic strengths and limitations. Evidence across reviewed studies suggests that this approach is particularly well-suited to analyzing crossing decision determinants under specific traffic conditions. Limitations of this approach include reduced generalizability to contexts with substantially different traffic signal configurations. Across these 3 methodological families, a clear tension emerges between interpretive richness and predictive performance. Structural models — including discrete choice and social force approaches — provide parameters directly interpretable for traffic engineering and policy purposes, whereas data-driven approaches achieve higher predictive accuracy on specific benchmark datasets but lack comparable policy utility. No single paradigm has been validated across the full range of pedestrian populations, traffic environments, and signal configurations represented in this corpus.

\section{External Interface Design for Autonomous Vehicles}
This section synthesizes findings from 3 papers examining Pedestrian crossing decision at crosswalk with approaching FAV, Feature preference evaluation, Receptivity assessment before/after exposure using Virtual Reality Experiment / Human-subjects laboratory experiment, Binomial logistic regression / Statistical modeling, Random parameter logit model (binary logit regression with normally distributed random parameters). The comparative evidence matrix structures claims, metrics, and limitations to support structured analysis. A key observation is that Virtual Reality Experiment / Human-subjects laboratory experiment and Binomial logistic regression / Statistical modeling approaches are the dominant methodological families in this corpus, though significant heterogeneity in reporting practices limits direct cross-study comparison.

\section{Factors Influencing Pedestrian Red-Light Running Behavior}
This section synthesizes findings from 2 papers examining Analyze pedestrian red light running decision-making, measure effects of personal factors (gender, age, companion presence), environmental factors (signal timing, traffic volume), and social influences (behavior of other pedestrians), Predict probability of pedestrian red-light running behavior using Random parameter logit model (binary logit regression with normally distributed random parameters), Binomial logistic regression / Statistical modeling. The comparative evidence matrix structures claims, metrics, and limitations to support structured analysis. A key observation is that Random parameter logit model (binary logit regression with normally distributed random parameters) and Binomial logistic regression / Statistical modeling approaches are the dominant methodological families in this corpus, though significant heterogeneity in reporting practices limits direct cross-study comparison.

\section{Cross-Study Comparison and Claim Synthesis}
Reported metrics in the reviewed literature span multiple dimensions of crossing behavior and traffic system performance, including: Odds ratios, coefficient estimates, z-values, AIC (Akaike Information Criterion) for model comparison, Score test p-values; Odds ratios; Model coefficients; 95% confidence intervals, Feature ratings (1-7 scale); Crossing time (seconds); Waiting time before crossing; Pedestrian receptivity scores; Personal innovativeness scores; PBQ subscales: violation, error, lapse, aggression, positive behaviors. Performance reporting is most complete for Analyze pedestrian red light running decision-making tasks. Metric reporting practices vary considerably across papers in two key respects. First, normalization approaches differ: some studies express outcomes as normalized rates per unit time or per pedestrian, while others report raw counts. Second, the temporal resolution of reported metrics ranges from instantaneous (e.g., crossing initiation delay at a single phase) to aggregate (e.g., violation frequency over a multi-hour observation window). These differences make direct cross-study comparison difficult without explicit standardization. Addressing this heterogeneity will require consensus on metric definitions and reporting standards within the pedestrian behavior modeling community. The absence of such standards currently forces reviewers to rely on narrative synthesis rather than formal meta-analysis, limiting the precision of conclusions that can be drawn.

\section{Research Gaps and Future Directions}
The comparative analysis of 2 papers reveals 1 persistent gaps that current approaches have not adequately resolved. These gaps are documented through verified evidence from the review matrix and are categorized by severity below. Gap [b887fd39] (medium severity): There are unresolved claim differences across papers under related tasks, suggesting a need for clearer comparison criteria.. This gap limits the ability to draw robust conclusions from the existing literature. These gaps collectively indicate that the field lacks the empirical foundation needed to support confident generalization of crossing behavior models across contexts, populations, and technologies. Prioritizing the most severe gaps in future empirical work will be essential for advancing both scientific understanding and the practical applicability of crossing behavior models.

\end{document}
% BIB START
@article{05_2018_deb_trf,
  author={Shuchisnigdha Deb; Lesley J. Strawderman; Daniel W. Carruth},
  title={Investigating pedestrian suggestions for external features on fully autonomous vehicles: A virtual reality experiment},
  year={2018},
  journal={Transportation Research Part F}
}

@article{12_2016_zhang_aap,
  author={Weihua Zhang; Kun Wang; Lei Wang; Zhongxiang Feng; Yingjie Du},
  title={Exploring factors affecting pedestrians' red-light running behaviors at intersections in China},
  year={2016},
  journal={Accident Analysis and Prevention}
}

@article{13_2020_chen_trf,
  author={Dianchen Zhu; N.N. Sze; Lu Bai},
  title={Roles of personal and environmental factors in the red light running propensity of pedestrian: Case study at the urban crosswalks},
  year={2020},
  journal={Transportation Research Part F}
}